\documentclass[11pt,a4paper]{article}
\usepackage[T1]{fontenc}
\usepackage[margin=1in]{geometry}
\usepackage{amsmath,amssymb,amsthm,mathtools}
\usepackage{booktabs}
\usepackage{xcolor}
\usepackage[colorlinks=true,linkcolor=blue!60!black,citecolor=blue!60!black,
            urlcolor=blue!60!black]{hyperref}
\usepackage{microtype}
\usepackage{enumitem}
\usepackage[most]{tcolorbox}

\newcommand{\lean}[1]{\texttt{\small #1}}
\newtcolorbox{keybox}[1][]{colback=blue!3,colframe=blue!45!black,
  boxrule=0.5pt,left=4pt,right=4pt,top=3pt,bottom=3pt,#1}
\newtcolorbox{failbox}[1][]{colback=red!3,colframe=red!45!black,
  boxrule=0.5pt,left=4pt,right=4pt,top=3pt,bottom=3pt,#1}
\newtheorem{theorem}{Theorem}

\title{\textbf{Duality is a structure; the sector is the cause}\\
\large T-duality of the compact boson stated on its sectors and checked by two kernels,\\
the measured invariant it predicts, and a criterion for when a self-dual point is physical}

\author{Xavier Callens\\
\small SocrateAI-Scientific-QuantumFluids\\
\small \url{https://github.com/xaviercallens/SocrateAI-Scientific-QuantumFluids}}
\date{September 2026}

\begin{document}
\maketitle

\begin{abstract}
\noindent
An involution $R\leftrightarrow\alpha'/R$ on a length scale was, in earlier work of this programme, offered a
physical home in T-duality, in K3 mirror symmetry, and in the modular symmetry of the quantum Hall system. Three
rounds of pre-registered numerics and twenty-eight machine-checked theorems later, we can say where such a
duality lives and what it does. We state T-duality on the sectors $(n,w)\in\mathbb Z^2$ of the compact boson and
prove, in Lean~4 against Mathlib with two independent kernels, that (i) it is a reindexing of the sector lattice
that leaves the sector sum invariant with no convergence hypothesis, (ii) the product of the electric and magnetic
scaling dimensions is radius-independent, $\Delta_e\Delta_m=\tfrac14$, which in the superfluid dictionary is
$\eta\,n_s\lambda_T^2=1$ -- the relation this programme measured to $12$--$27\%$ on eleven energies in two boxes --
and (iii) the self-dual radius $\sqrt2$ and the radius $2\sqrt2$ at which the vortex becomes marginal are
different points: the Berezinskii--Kosterlitz--Thouless transition is the threshold of one sector, not the fixed
point of the exchange. From these three facts and the literature on Ising, string-gas and superconductor--insulator
dualities we extract a criterion: a self-dual point pins a transition only when the dual sectors are simultaneously
critical on a discrete set of fixed points; on a line of fixed points it is a point like any other; when both
sectors are populated it can be a preferred equilibrium. In every case what acts is a sector -- a conserved label
that crosses a threshold, proliferates, annihilates, or is imposed -- and the duality relabels it. The duality is
exact, verified, and not a cause. No physics here is new; the mathematics is elementary; what is new is that the
distinction is stated so that a kernel can check it and an experiment can test it.
\end{abstract}

\section{Three questions one involution raised}

The dual-length involution $k\mapsto k_s^2/k$ of~\cite{Callens2026Lean} has the form of T-duality. It was pointed
out that the same involution is the Fricke element of $\Gamma_0(N)^+$, the mirror involution of degree-$2N$
lattice-polarised K3 surfaces~\cite{Dolgachev1996}, and that $\Gamma_0(2)$ acts on the quantum Hall conductivity
plane~\cite{LutkenRoss1992}. Each identification was taken to its gate~\cite{Callens2026Wasserstein,
Callens2026Causal}: the Fricke matrix normalises Mathlib's $\Gamma_0(n)$ and its axis action is the dual-length map,
but no $\Gamma_0(n)$ acts on a wavenumber (\lean{Fricke.lean}); the Fricke element fixes the Hall critical point's
orbit but sends every plateau off the odd-denominator class (\lean{QHFricke.lean}); and the physical BKT point is not
the self-dual radius. The residue of these three verdicts is one question: if a duality is exact and verified, and
yet is not where the physics happens, what is it, and what is?

\section{T-duality on the sector lattice, machine-checked}

Let the compact boson have radius $R$ in the normalisation where the self-dual radius is $\sqrt2$. The primary of
momentum $n$ and winding $w$ has $h=\tfrac12(n/R+wR/2)^2$, $\bar h=\tfrac12(n/R-wR/2)^2$, hence
\[
\Delta(n,w;R)=n^2/R^2+w^2R^2/4,\qquad s=h-\bar h=nw .
\]
The electric (spin-wave, $e^{i\theta}$) operator is $(1,0)$; the magnetic (vortex) operator is $(0,1)$. The
following are theorems of \lean{CompactBoson.lean} (twelve statements, axiom footprint $\{\mathrm{propext},
\mathrm{Classical.choice},\mathrm{Quot.sound}\}$, accepted by the Lean kernel and by \lean{nanoda} through
Comparator; three negative controls -- product $\tfrac12$, BKT at $\sqrt2$, duality with $R\mapsto1/R$ -- fail).

\begin{theorem}[duality as reindexing]
$\Delta(n,w;R)=\Delta(w,n;2/R)$ and $s(n,w)=s(w,n)$. Consequently, for every weight $F$,
$\sum_{(n,w)\in\mathbb Z^2}F(\Delta(n,w;R),s)=\sum_{(n,w)\in\mathbb Z^2}F(\Delta(n,w;2/R),s)$.
\end{theorem}
The second statement is \lean{Equiv.tsum\_eq} applied to the swap of $\mathbb Z^2$: reindexing a sum by a
bijection is an identity of \lean{tsum} whether or not the series converges. This is what the duality is -- a
bijection of the sector lattice that preserves the spectrum. Savit's review~\cite{Savit1980} gives the general
form: a duality is a Fourier transform on the group of sectors, electric charges in $G$ against magnetic charges
in $\hat G$, with the Lagrangian sum over one family turned into the Hamiltonian sum over the dual family by Poisson
summation.

\begin{theorem}[the invariant]
$\Delta(1,0;R)\,\Delta(0,1;R)=\tfrac14$ for every $R\neq0$.
\end{theorem}
In the XY / two-dimensional-superfluid dictionary $\Delta(1,0)=\eta/2$ and $\Delta(0,1)=n_s\lambda_T^2/2$, so the
invariant reads $\eta\,n_s\lambda_T^2=1$. That is the relation measured in~\cite{Callens2026Causal}: $1.047$ and
$1.072$ on two equilibrated energies in round~2, $1.12$--$1.22$ on the six energies of the heating ladder,
$1.14$--$1.27$ on the five energies of the $L=32$ box -- with the finite-size excess predicted by
Hasenbusch's logarithmic corrections~\cite{Hasenbusch2005}. \emph{The numerical verification of the duality in
this programme is the measurement of this product.}

\begin{theorem}[self-dual is not marginal]
For $R>0$: $\Delta(1,0)=\Delta(0,1)$ iff $R=\sqrt2$, where both equal $\tfrac12$; $\Delta(0,1)=2$ iff
$R=2\sqrt2$, where $\Delta(1,0)=\tfrac18$, i.e.\ $\eta=\tfrac14$; and $2\sqrt2\neq\sqrt2$.
\end{theorem}
The transition is where the vortex operator becomes marginal -- Nelson and Kosterlitz's universal jump
$n_s\lambda_T^2=4$~\cite{NelsonKosterlitz1977}, located in our own runs at $T_{\rm BKT}=0.821$ with
$\eta=0.31$ there -- and that is a condition on one sector. The exchange symmetry has its fixed point elsewhere,
and nothing happens there.

\section{When is a self-dual point physical?}

The compact boson is the case with a \emph{line} of fixed points, one for each $R$. Three other cases complete the
criterion.

\paragraph{Ising: yes.} Kramers--Wannier self-duality~\cite{KramersWannier1941} fixes $T_c$ at $\sinh2K=1$ because
the order and disorder operators~\cite{FradkinSusskind1978} exchange under the duality and $\mathbb Z_2$ has a single
symmetric point: the dual sectors are simultaneously critical and the fixed-point structure is discrete.

\paragraph{String gas cosmology: a preferred equilibrium.} In the Brandenberger--Vafa
scenario~\cite{BrandenbergerVafa1989,Brandenberger2008} ``the symmetry we make use of is T-duality, and the new
degrees of freedom are the string oscillatory modes and the string winding modes''. Momentum modes, with energy
$n/R$, resist contraction; winding modes, with energy $wR$, resist expansion; when both are populated the radion is
stabilised at the self-dual radius as the minimum of an effective potential. The self-dual point is physical there
because both dual sectors are present and compete. And the exit from the Hagedorn phase -- the event that makes
three dimensions large -- is the \emph{annihilation of winding modes}: a sector event, with T-duality as the
symmetry of the spectrum in which it is described.

\paragraph{Superconductor--insulator transition: a constraint.} Boson--vortex duality~\cite{DasguptaHalperin1981,
FisherLee1989} predicts, if the transition is self-dual, a universal resistance $h/4e^2$~\cite{Fisher1990,
ChaFisherGirvin1991}; measured values scatter around it. The cause of the transition is the proliferation of
vortices on one side and of charges on the other; self-duality constrains the critical point only if the fixed
point happens to be self-dual, which is a property of the fixed point, not of the duality.

\paragraph{Quantum Hall: organisation, not cause.} $\Gamma_0(2)$ organises the plateau diagram and the position
of the critical points~\cite{LutkenRoss1992}; the cause of the quantisation is Laughlin's gauge
argument~\cite{Laughlin1981} on the Chern sector; and the Fricke element, which fixes the critical orbit, is not a
symmetry of the plateaux (\lean{QHFricke.lean}).

\begin{keybox}
\textbf{Criterion.} A self-dual point pins a transition iff the two dual sectors are simultaneously critical and
the fixed-point set is discrete. On a line of fixed points the self-dual point is a point like any other and the
transition sits where one sector crosses its own threshold. When both dual sectors are populated and compete
energetically, the self-dual point is a preferred equilibrium. In each case the duality relabels the sectors and
the sector acts.
\end{keybox}

\section{What acts: the sector}

The same triad recurs in every case above: a group of sectors ($\mathbb Z$ for $U(1)$ windings, $\mathbb Z_2$ for
Ising, Chern numbers, string windings); a duality that is a Fourier or reindexing map between $G$-sectors and
$\hat G$-sectors preserving the spectrum; and physics that happens when a sector crosses a threshold, proliferates,
annihilates, or is imposed. The causal role of the sector is what the rest of this programme
established~\cite{Callens2026Causal}: it is conserved by every continuous evolution without a phase slip
(\lean{TopologicalProtection}), additive in scale so that a bound pair is invisible from outside
(\lean{ScaleResolvedWinding}), read correctly by the sampled loop sum as the continuum degree
(\lean{ContinuumWinding}), it carries its own temperature (\lean{SectorTemperature}), and in an energy-matched
intervention it is a difference-maker: the same energy delivered as topology lowered the condensate by
$0.49$--$0.73$ where phonons lowered it by at most $0.035$, with the injected sector reading five to eight times
the bath's temperature on our own calibrated thermometer. That is why one involution appears in helium, Hall bars,
K3 moduli and string gases, and why it explains none of them.

\section{What this paper does not claim}
\begin{failbox}
\begin{itemize}[leftmargin=*]
\item Novelty in the mathematics: the three theorems are elementary; their value is that the distinction between
  the duality and the sector is stated so that a kernel checks it.
\item Novelty in the physics: Kramers--Wannier, Nelson--Kosterlitz, Brandenberger--Vafa, Fisher and Lütken--Ross
  are cited, not extended.
\item Any identification of a superfluid with a K3 surface or a string background. The K3 mirror involution is the
  same map on a different space; the criterion above is what the shared map does and does not license.
\item That every duality in physics fits the triad: the cases examined do; the general statement is a reading.
\end{itemize}
\end{failbox}

\section*{Addendum (2026-09-26): the criterion sharpened by the non-invertible-symmetry literature}
\addcontentsline{toc}{section}{Addendum}
A literature review after publication located the modern machinery that makes \S3's criterion precise rather
than descriptive. Aasen, Mong and Fendley state the core distinction directly: a self-duality is implemented by
an operator that ``is not a symmetry in the traditional sense... not unitary or even invertible''
\cite{AasenMongFendley2016}. Choi, C\'ordova, Hsin, Lam and Shao exhibit the Poisson resummation that turns this
into an explicit lattice statement, and prove that at some self-dual points a mixed 't~Hooft-like anomaly
\emph{forbids} a trivially gapped phase -- an anomaly argument, not the duality map, is what does the forbidding
\cite{ChoiCordovaHsinLamShao2022}. Kaidi, Ohmori and Zheng give the clean mechanism for our own already-published
observation: at $SO(3)$ Yang--Mills $\theta=\pm\pi$, the self-dual point sits \emph{between two physically
inequivalent gapped phases} (a $\mathbb Z_2$ TQFT and the trivial phase), which is why it is forced to be a
transition \cite{KaidiOhmoriZheng2022}. Shao's lecture notes give the general rule, in one sentence: away from a
critical point Kramers--Wannier duality ``is not a duality relating two equivalent descriptions of a single
system... it is a map from the high-temperature phase to the low-temperature phase,'' and at the compact boson's
self-dual radius $\tau=i$ ``the duality defect reduces to an invertible $\mathbb Z_4$ symmetry'' \cite{Shao2023} --
exactly why $R=\sqrt2$ (a point on a continuous moduli space with enhanced symmetry) is structurally unlike the
Ising or $SO(3)$-Yang--Mills self-dual points (between two gapped, inequivalent phases), and why neither is like
the vortex-marginality radius $2\sqrt2$, which is a threshold on one sector alone. One further correction:
Hayashi and Tanizaki show that the Cardy--Rabinovici $SL(2,\mathbb Z)$ ``self-duality'' of the 2D Coulomb
gas/$\mathbb Z_N$ clock model is not naive -- it holds only after gauging the electric one-form symmetry with a
discrete $\theta$-term -- and that the resulting self-dual point, where three first-order transition lines meet,
is again forced non-trivial by a separate mixed gravitational anomaly, not by the duality symmetry itself
\cite{HayashiTanizaki2022}. None of this changes a theorem of this paper; it replaces the numerical observation
``$\sqrt2\neq2\sqrt2$'' with an explicit criterion (inequivalent sectors on the two sides $\Rightarrow$ forced
transition; a continuous moduli space with enhanced symmetry $\Rightarrow$ no transition) that a future module
could formalise as a finite-group gauging statement.

\section*{Addendum (2026-09-26, part 2): two more cases, same criterion}
\addcontentsline{toc}{section}{Addendum 2}
Two further cross-domain instances, found by literature review and independently verified, apply the criterion of
\S3 without changing it.

\paragraph{RCFT/Verlinde: the sharp case, and the same verdict as BKT.} The Verlinde formula~\cite{Verlinde1988},
proved from Moore--Seiberg's polynomial equations~\cite{MooreSeiberg1988,MooreSeiberg1989}, makes the modular
$S$-matrix acting on the primaries (the sectors) of a rational conformal field theory a duality satisfying
$S^2=C=(ST)^3$ exactly, for $S$ unitary and symmetric~\cite{Fuchs1993} -- sharper than this programme's own
$\mathbb Z_N$ case, where Mathlib's unnormalised \lean{ZMod.dft} gives $\mathrm{dft}\circ\mathrm{dft}=N\cdot
\Phi(-\cdot)$, not the identity (\lean{SectorDuality.lean}). $S$ is a literal finite-abelian-group Fourier
transform only for simple-current (pointed) fusion rings -- e.g.\ $\widehat{su}(2)_1$, whose two primaries fuse as
$\mathbb Z_2$ and whose $S$ is exactly the (unnormalised) Hadamard matrix
$\left(\begin{smallmatrix}1&1\\1&-1\end{smallmatrix}\right)$ -- and here $S^2=C=\mathbb 1$ with no extra factor
once the missing $1/\sqrt2$ is restored (\lean{RCFTDuality.su2Level1\_S\_sq}), the clean identity the unnormalised
$\mathbb Z_N$ statement only approximates up to the decoupled factor $|G|=2$. For a generic (non-abelian) fusion
category $S$ is instead the Kac--Peterson matrix and no group Fourier transform underlies it at all; this module
proves only the $2\times2$ integer-matrix identity, not that it is the Kac--Peterson matrix, which is affine-Lie-
algebra representation theory supplied here as physics input. The natural self-dual point of the worked example --
the free-boson radius $R=1$ where $\hat u(1)$ enhances to $\widehat{su}(2)_1$ under T-duality $R\leftrightarrow1/R$
-- sits on a \emph{continuous} line of RCFT fixed points, $1\le R\le\sqrt2$, separating three gapped
phases~\cite{PaceChatterjeeShao2024}, not on a discrete pair of inequivalent sectors: by this programme's own
criterion it is a point like any other, not a forced transition -- the Ising verdict does not transfer, the BKT
verdict does.

\paragraph{Particle-vortex/level-rank in the fractional quantum Hall effect: a second, orthogonal duality.} At
$\nu=\tfrac12$ the Halperin--Lee--Read composite fermion is conjectured to be a massless Dirac fermion, exactly
self-conjugate under the particle-hole (CT) involution that exchanges the Jain-sequence pair
$n/(2n+1)\leftrightarrow(n+1)/(2n+1)$~\cite{Son2015}; Seiberg--Senthil--Wang--Witten embed this in a fermionic
particle-vortex duality web~\cite{SeibergSenthilWangWitten2016} whose rigorously-established topological limit is
level-rank duality, $U(N)_K\leftrightarrow U(K)_N$~\cite{HsinSeiberg2016,NaculichSchnitzer2007}. This is a
different duality from $\Gamma_0(2)$ on the conductivity plane, acting on a different sector: the finite set of
Chern--Simons/WZW primaries, labelled by Young diagrams in an $N\times K$ box, relabelled by transpose to a
$K\times N$ box -- an operation Mathlib already formalises on \lean{YoungDiagram} as an order isomorphism
(\lean{transposeOrderIso}), restricted here to two finite boxes (\lean{LevelRankDuality.levelRankRelabeling}).
Filling $\nu=\tfrac12$ is not itself a transition between two gapped, inequivalent sectors -- that role belongs to
the Jain pairs on either side of it -- but the unique fixed point of the particle-hole involution on the
continuous filling $\nu$, in the same slot as the compact boson's self-dual radius: an enhanced-symmetry point,
not a cause. What self-duality there fixes is a number, not a phase boundary -- the composite fermion's Berry
phase is forced to be exactly $\pi$ and $\sigma^{CF}_{xy}=-\tfrac12$ exactly~\cite{Son2015} -- so this instance
sits with the superconductor--insulator transition in the \emph{constraint} slot, not with $\Gamma_0(2)$'s
\emph{organisation} slot. Level-rank duality itself, wherever it appears in the non-Abelian FQH hierarchy, is
unambiguously organisation-not-cause: a proved relabelling of a finite anyon label set that decides no physics
and causes no transition. Neither Lean module proves the physical content of its own duality (S,T-matrix or
Verlinde-formula preservation for level-rank; that the Hadamard matrix is the Kac--Peterson matrix for RCFT) --
in both cases the module certifies a decidable combinatorial or linear-algebraic skeleton, and the identification
with the named physical duality is supplied as literature-sourced prose, exactly as \lean{SectorDuality.lean} and
\lean{QHFricke.lean} already do.

\section*{Addendum (2026-09-26, part 3): a ninth case, and one explicitly declined}
\addcontentsline{toc}{section}{Addendum 3}
A further literature pass, following the same criterion, adds one more continuous-window case and explicitly
excludes a candidate that turned out, on reading, not to be a duality at all.

\paragraph{Seiberg duality: a third continuous window, same verdict.} Electric SU($N_c$) SQCD with $N_f$ flavors
and its magnetic dual SU($N_f-N_c$)~\cite{Seiberg1994} flow to the same interacting fixed point throughout the
conformal window $\tfrac32 N_c<N_f<3N_c$~\cite{IntriligatorSeiberg1996} -- ``completely equivalent at long
distance,'' not two phases exchanged by the duality. The sector here is the integer $N_c$ itself, relabelled by
$N_c\mapsto N_f-N_c$; the point $N_c=N_f/2$, where the two gauge groups coincide, is a direct algebraic fact
about this map, not a distinguished physical locus -- no transition occurs there, the whole window being one
phase. Same verdict as the compact boson and the RCFT self-dual radius: a point on a continuum, not a cause. No
Lean statement is proposed: the only formalisable content of the map $N_c\mapsto N_f-N_c$ is a one-line integer
involution, and the actual physical duality (holomorphy, anomaly matching, RG-flow equivalence) has no
combinatorial skeleton analogous to \lean{YoungDiagram.transpose} or \lean{ZMod.dft}; unlike Kramers--Wannier or
T-duality, the two Lagrangians are not even operator images of each other (the magnetic side carries an extra
gauge-singlet meson with its own superpotential).

\paragraph{Topological superconductors: not a duality.} Reading Kitaev's classification of unpaired Majorana
modes~\cite{Kitaev2001}, the Altland--Zirnbauer symmetry classes~\cite{AltlandZirnbauer1997}, and
Ryu--Schnyder--Furusaki--Ludwig's tenfold-way classification~\cite{RyuSchnyderFurusakiLudwig2010} directly finds
no duality organising this material: it is a classification of gapped free-fermion phases by Clifford-algebra
K-theory and Bott periodicity, related by dimensional reduction, not an equivalence between two theories. It has
a sector (a $\mathbb Z$ or $\mathbb Z_2$ topological invariant) but no map that relabels it against a dual
description -- the word ``duality'' organises none of these three sources. This is therefore \emph{excluded},
not merely deferred: forcing this programme's duality/sector vocabulary onto a classification would be the
formalising-an-analogy failure this programme's own discipline exists to catch. One legitimate but
pre-existing link is worth recording as a footnote rather than a new case: the 1D Kitaev chain (class BDI) is
related by the Jordan--Wigner transform to the transverse-field Ising model, so its topological/trivial
transition is, in fermionic language, the same Kramers--Wannier self-duality already the Ising case above --
not a ninth instance, a relabelling of the first one.

\section*{Addendum (2026-09-26, part 4): a tenth case, checked almost entirely by Mathlib itself}
\addcontentsline{toc}{section}{Addendum 4}
A further round of literature review, looking specifically for new cases beyond every one already covered here,
found one strong addition and two candidates that, on reading, do not fit the criterion's shape at all.

\paragraph{Montonen--Olive $S$-duality: order-4 and order-6 points on a continuous coupling, not a cause.}
$\mathcal N=4$ super Yang--Mills has an exactly marginal complexified coupling $\tau=\theta/2\pi+4\pi i/g^2$
ranging over the upper half-plane; $S$-duality~\cite{MontonenOlive1977} conjectures a quantum symmetry
$S:\tau\mapsto-1/\tau$, exchanging the gauge group with its Langlands dual and electric with magnetic charges,
which together with the classical shift $T:\tau\mapsto\tau+1$ generates $\mathrm{SL}(2,\mathbb
Z)$~\cite{KapustinWitten2007}. Mathlib already proves the two facts this needs, with no new theorem from this
project: \lean{ModularGroup.S\_mul\_S\_eq} ($S^2=-1$, so $S$ has order four) and
\lean{ModularGroup.stabilizer\_I} (the stabiliser of $\tau=i$ in $\mathrm{SL}(2,\mathbb Z)$ is exactly
$\{\pm1,\pm S\}$) -- the literal statement that $\tau=i$ has enhanced discrete symmetry, and, by
\lean{ModularGroup.stabilizer\_rho}, so does $\rho=e^{i\pi/3}$ (order six). As with every continuous modulus in
this programme, this changes nothing about the physics: $\mathcal N=4$ SYM is superconformal for every $\tau$,
there is no phase transition anywhere on the upper half-plane, and $\tau=i$ is a point like any other except for
the extra unbroken $\mathbb Z_4\subset\mathrm{SL}(2,\mathbb Z)$ -- the BKT/RCFT/Seiberg verdict, not the Ising
one. What this does NOT resolve: at $\tau=i$, $S$ maps the theory to itself only after also exchanging discrete
global variants of the gauge group (e.g.\ $SU(N)$ with $PSU(N)$, physically distinct theories with different
line-operator spectra); whether that additional exchange sharpens the verdict is a further question this round
does not settle.

\paragraph{Two candidates excluded as structural mismatches, not merely deferred.} Three-dimensional Ising/Wegner
duality~\cite{Wegner1971} maps the Ising spin model to a $\mathbb Z_2$ \emph{lattice gauge theory} on the dual
lattice -- a genuinely different theory with local gauge symmetry and link degrees of freedom, not a self-map;
self-duality closing on the same model is a two-dimensional coincidence, already this programme's Ising case.
There is no fixed point to ask a self-dual-point question about in $d=3$: the criterion does not apply
structurally, not merely give a negative answer. AdS/CFT~\cite{Maldacena1998} is a strong--weak
\emph{equivalence} between two descriptions of one theory, holding at every 't~Hooft coupling; there is no
conjectured involution on the coupling with a fixed locus of enhanced symmetry the way $S$-duality has one at
$\tau=i$ -- it relates two variables describing a single moduli-free correspondence, not a relabelling of
sectors on a moduli space with a distinguished point, so the self-dual-point question this programme asks does
not apply to it either.

\section*{Addendum (2026-09-26, part 5): the T-duality thread's furthest extension, and one question left
open on purpose}
\addcontentsline{toc}{section}{Addendum 5}
The previous round flagged Calabi--Yau mirror symmetry beyond K3$\times T^2$ as researched not at all, only
guessed at. This round replaces the guess with a real search, and reports a genuinely new kind of outcome: not
a fourth ``continuous modulus, not forced'' verdict, and not an exclusion, but an open question stated as such.

\paragraph{Calabi--Yau mirror symmetry: the T-duality thread's furthest extension, and one question left open on
purpose.} Strominger, Yau and Zaslow show that mirror symmetry between a Calabi--Yau threefold $X$ and its
mirror $Y$ is literally T-duality on the three circles of a special Lagrangian $T^3$ fibration of
$X$~\cite{StromingerYauZaslow1996} -- the sector is the pair of continuous moduli spaces (K\"ahler moduli of
$X$, complex-structure moduli of $Y$), and this is the furthest geometric extension of this programme's own
T-duality thread: one circle for the compact boson, the Fricke involution on K3$\times T^2$, now a full $T^3$.
Unlike every case above, this is not manifestly a self-map of one moduli space -- mirror symmetry generically
relates \emph{two different} manifolds -- so the question this programme otherwise asks (is a self-dual point
forced) only becomes literal in the special case $X\cong Y$. We did not find, and do not assert, an established
example of a physically-distinguished self-mirror threefold playing a role analogous to $\tau=i$ or $R=\sqrt2$;
asserting one without a citable source would be exactly the unverified claim this programme's discipline exists
to prevent. What is established, and worth recording instead, is that mirror symmetry can fail outright rather
than close on a point: a rigid Calabi--Yau ($h^{2,1}=0$) has no ordinary smooth mirror, since the mirror would
need $h^{1,1}=0$, impossible for any compact K\"ahler manifold~\cite{CandelasDerrickParkes1993}. No Lean
statement is proposed: Mathlib has no Calabi--Yau, special-Lagrangian, or mirror-symmetry content to build on
(confirmed by direct search of the pinned checkout).

\section*{Addendum (2026-09-27): a third forced case, and one existing case's own self-dual point examined}
\addcontentsline{toc}{section}{Addendum 6}
The ``forced'' verdict has so far been thin -- Ising and $SO(3)$ Yang--Mills the only two instances
among nine other cases. This round adds a third, rigorously, and examines whether an already-integrated
case has a genuine self-dual point of its own.

\paragraph{Potts/random-cluster self-duality: Ising's family, made rigorous for every $q$.} The
$q$-state Potts model's random-cluster (Fortuin--Kasteleyn) representation has a duality, generalising
Kramers--Wannier, relating edge-weight $p$ on a planar graph to $p^\star(p;q)$ on its dual, with a
unique self-dual point $p_{\rm sd}(q)=\sqrt q/(1+\sqrt q)$. Beffara and Duminil-Copin prove, for every
$q\ge1$, that this point \emph{is} the critical point~\cite{BeffaraDuminilCopin2012} -- a rigorous,
all-$q$ generalisation of Ising's own $q=2$ case, not a physicist's heuristic. What self-duality does
not determine is the transition's character: it is continuous for $1\le q\le4$~\cite{DuminilCopinSidoraviciusTassion2017}
and discontinuous for $q>4$~\cite{DuminilCopinGagnebinHarelManolescuTassion2021} -- two sectors that
meet smoothly in one regime and jump in the other, at the exact same self-dual point either way. By
this programme's criterion every $q\ge1$ case is \emph{forced}: self-duality pins the transition's
location exactly regardless of its order, a fact now proved rather than assumed. As with Ising itself,
no Lean module is proposed: the checkable content is an elementary algebraic fixed point, and this
programme has never formalised Ising's own duality as a project module either.

\paragraph{Level-rank duality at $N=K$: the labels close up; the physics does not.} The Young-diagram
box-transpose bijection already proved in this programme (\lean{LevelRankDuality.levelRankRelabeling})
holds for every $a,b$ including $a=b$, so $N{\times}N$-box diagrams are trivially permuted by
transpose -- no new theorem is needed to say so. This combinatorial closure does not correspond to a
physical self-duality: Hsin-Seiberg's own $N=K$ computation shows $SU(N)_N$ relates not to itself but
to its level-reversal $SU(N)_{-N}$ tensored with an extra decoupled factor~\cite{HsinSeiberg2016}, and
neither primary source treats $N=K$ as a distinguished locus in this programme's sense. As with
AdS/CFT and three-dimensional Ising duality, $N=K$ poses no genuine self-dual-point question the
literature answers: asserting one from the label combinatorics alone, when the physics does not
support it, would be exactly the failure mode of formalising an analogy rather than a dependency.

\bibliographystyle{unsrt}
\bibliography{refs_duality_sector}
\end{document}
